\documentclass[11pt,a4paper]{article}

\usepackage[margin=1in]{geometry}
\usepackage{amsmath, amssymb, amsthm}
\usepackage{mathtools}
\usepackage{bm}
\usepackage{hyperref}
\usepackage{enumitem}
\usepackage{microtype}
\usepackage[T1]{fontenc}
\usepackage{lmodern}
\usepackage[dvipsnames]{xcolor}

\hypersetup{
  colorlinks=true,
  linkcolor=blue!50!black,
  urlcolor=blue!50!black,
  citecolor=blue!50!black
}

\newtheorem{theorem}{Theorem}[section]
\newtheorem*{theorem*}{Theorem}

\title{Relativistic Quantum Mechanics as a Forced Structural Consequence of Event-Density Primitives}

\author{Allen Proxmire \and Copilot (AI collaborator)}

\date{April 2026}

\begin{document}

\maketitle

\begin{abstract}
We show that relativistic quantum mechanics --- including the Klein--Gordon equation, the spin--statistics correspondence, the Clifford-algebra spinor frame, and the Dirac equation --- emerges as a \emph{forced structural consequence} of the Event-Density (ED) primitive stack, rather than as a set of independent postulates. Building on the Phase-1 derivation of non-relativistic quantum mechanics from the participation measure $P_K = \sqrt{b_K} \cdot e^{i\pi_K}$, we extend the program to the relativistic regime via three connected structural theorems. \textbf{Theorem R1} (spin--statistics) establishes $\eta = (-1)^{2s}$ at primitive level via the chain: exchange dichotomy $\to \pi_1(Q_2) = \mathbb{Z}_2 \to$ exchange-rotation generator identity $\to$ Clifford frame $\to$ minimal-bilinear-pairing closure. \textbf{Theorem R2} (Cl(3,1) uniqueness) shows that the real Clifford algebra with anticommutator $\{\gamma^\mu, \gamma^\nu\} = 2\eta^{\mu\nu}$ is the unique algebraic structure compatible with Lorentz covariance, half-integer representation content, and the rotational double cover; pure commutation and pure anticommutation alternatives are ruled out. \textbf{Theorem R3} (Dirac emergence) derives the Dirac equation $(i\gamma^\mu \partial_\mu - mc/\hbar)\Psi = 0$ as the unique first-order Lorentz-covariant linear equation on the Cl(3,1) spinor module, with positive-definite conserved current $j^\mu = \bar\Psi\gamma^\mu\Psi$, gauge-covariant minimal coupling, and tree-level $g = 2$ from non-relativistic reduction. ED inherits dimensional constants $(\hbar, c, m, q)$; it derives the \emph{form} of relativistic quantum equations.
\end{abstract}

\section{Introduction}

The mathematical structure of quantum mechanics --- wavefunctions on Hilbert space, unitary evolution, the Born rule, exchange statistics, spinor representations of the Lorentz group, and the first-order Dirac equation --- is conventionally introduced as a layered set of independent postulates. Each layer (kinematic, dynamical, statistical, relativistic, gauge-theoretic) is justified primarily by its empirical success rather than by derivation from a more fundamental ontology.

The Event-Density (ED) program offers an alternative: to build quantum mechanics as a \emph{forced structural consequence} of a small primitive stack governing micro-events, chains, bandwidth content, and commitment dynamics on a 3+1-dimensional event manifold. Phase-1 of this program~\cite{phase1} established that non-relativistic single-particle quantum mechanics --- the Schr\"odinger equation, the Born rule, Bell--CHSH violations saturating the Tsirelson bound, and the Heisenberg uncertainty relation --- emerges from the participation measure
\begin{equation}
P_K(x, t) = \sqrt{b_K(x,t)} \cdot e^{i\pi_K(x,t)},
\end{equation}
combined with the four-band bandwidth decomposition (Primitive~04) and commitment dynamics (Primitives~10--11). The squared-amplitude exponent $b_K = |P_K|^2$ is \emph{definitional}, not postulated; the Born rule's $|\psi|^2$ is therefore not an independent axiom. The same exponent-2 structure threads through Madelung's $\sqrt{\rho}$, Schr\"odinger's kinetic term, sublinear bandwidth composition, and Heisenberg's $(\Delta x \cdot \Delta p)^2$, yielding ED's strongest positive structural unification.

Phase-2 extends the program. Arc~R, the subject of the present paper, derives \textbf{relativistic single-chain quantum mechanics} at primitive level. Three structural theorems are established. Arc~M~\cite{arcm} (chain-mass structure) and Arc~Q~\cite{arcq} (QFT extension) build on Arc~R; we summarise their relationship in Section~\ref{sec:discussion}.

The motivation for deriving rather than postulating is threefold. \emph{First}, structural derivation reduces the axiomatic surface area of physics: postulates that turn out to be theorems carry no independent empirical content and can be eliminated. \emph{Second}, structural derivation exposes the precise inputs each piece of quantum theory requires; this clarifies which features are inherited (numerical constants, specific particle content) versus structural (form of equations, classification of representations). \emph{Third}, structural derivation enables systematic extension: if non-relativistic quantum mechanics is a thin-participation limit of a richer primitive structure, the same primitives extend to relativistic and field-level regimes by relaxing the limit, rather than by introducing new postulates.

The Arc~R results consist of three theorems plus extensive supporting structural content:
\begin{itemize}
  \item \textbf{R1 (Spin--statistics):} $\eta = (-1)^{2s}$ derived from primitive-level exchange topology and frame structure;
  \item \textbf{R2 (Cl(3,1) uniqueness):} the real Clifford algebra with anticommutator $\{\gamma^\mu, \gamma^\nu\} = 2\eta^{\mu\nu}$ is the unique algebraic frame compatible with the Lorentz / spin / topology constraints;
  \item \textbf{R3 (Dirac emergence):} the Dirac equation is the unique first-order Lorentz-covariant equation on the Cl(3,1) spinor module, with conserved current and tree-level $g = 2$ as automatic consequences.
\end{itemize}

ED's primitive-level account does \emph{not} predict the numerical values of dimensional constants ($\hbar$, $c$) or rule-type-specific parameters ($m$, $q$); these are inherited via the Dimensional Atlas. The clean separation between \emph{form} (forced structurally) and \emph{value} (inherited empirically) is a load-bearing methodological principle that recurs throughout Arc~R and continues into Arc~M and Arc~Q.

\section{Background}

\subsection{The ED primitive stack}

Event Density is built on thirteen primitives. The primitives relevant to Arc~R are summarised below:
\begin{itemize}
  \item \textbf{Primitive 01 (micro-event):} discrete events on the 3+1-dimensional ambient manifold; substrate for all participation.
  \item \textbf{Primitive 02 (chain):} worldlines $\gamma_K$ of persistent participation.
  \item \textbf{Primitive 03 (participation):} act of a chain entering a micro-event with non-zero amplitude.
  \item \textbf{Primitive 04 (bandwidth):} four-band decomposition $b_K = b_K^{\mathrm{int}} + b_K^{\mathrm{adj}} + b_K^{\mathrm{env}} + b_K^{\mathrm{com}}$.
  \item \textbf{Primitive 06 (four-gradient):} $\partial_\mu$ on the event manifold; carries Lorentz covariance.
  \item \textbf{Primitive 07 (rule-type):} discrete classification by structural levers L1--L4.
  \item \textbf{Primitive 09 (polarity phase):} $\pi_K$ of the participation measure.
  \item \textbf{Primitive 10 (individuation):} threshold separating distinct chains.
  \item \textbf{Primitive 11 (commitment):} dynamical update producing micro-event events along $\gamma_K$.
  \item \textbf{Primitive 13 (relational timing):} proper time $\tau_K$ along $\gamma_K$.
\end{itemize}

\subsection{The participation measure}

The Phase-1 participation measure
\begin{equation}
P_K(x^\mu) = \sqrt{b_K(x^\mu)} \cdot e^{i\pi_K(x^\mu)} \in \mathbb{C}
\end{equation}
fuses Primitive~04 (bandwidth amplitude) with Primitive~09 (polarity phase). Stage~R.1 promotes this object to a Lorentz-covariant scalar field on the ambient 3+1D event manifold, with $x^\mu = (ct, \mathbf{x})$. Lorentz covariance is forced by Primitive~06.

\subsection{Rule-type ontology}

A rule-type $\tau$ is specified by four levers from Primitive~07:
\begin{itemize}
  \item \textbf{L1:} bandwidth partition pattern $w_\tau^X$ across the four bands;
  \item \textbf{L2:} internal-index content (Lorentz representation $(j_L, j_R)$ plus gauge index);
  \item \textbf{L3:} interface content (Fierz class $\Gamma_\tau$ or analogue);
  \item \textbf{L4:} statistics class (Case~P bosonic / Case~R fermionic).
\end{itemize}
Stage~R.2 of Arc~R produces the spin--statistics correspondence connecting L4 to L2.

\subsection{Commitment and individuation}

Primitive~11 supplies the dynamical update along chain worldlines: discrete events in which a chain ``commits'' --- depositing a definite outcome on the event manifold. Primitive~10 sets the threshold at which two same-type chains remain distinguishable; for Case-R rule-types, individuation forbids same-type coincidence, leading to antisymmetry and Pauli exclusion. The interplay of these two primitives is central to the spin--statistics theorem of Section~\ref{sec:r2}.

\section{Arc R Stage R.1 --- Scalar Relativistic Quantum Mechanics}

\subsection{Klein--Gordon equation}

The Stage~R.1 derivation begins with the Lorentz-covariant participation measure on the $(0,0)$ representation. For a free scalar rule-type, requiring (i) Lorentz-scalar form of the dynamical equation, (ii) translation invariance, and (iii) reduction to the non-relativistic Schr\"odinger limit yields the Klein--Gordon equation
\begin{equation}
\left( \Box + \frac{m^2 c^2}{\hbar^2} \right) \Psi = 0,
\label{eq:KG}
\end{equation}
where $\Box = \eta^{\mu\nu} \partial_\mu \partial_\nu$. The mass term $m^2 c^2/\hbar^2$ enters as the Casimir eigenvalue of the Poincar\'e algebra acting on the rule-type's participation module; numerical $m$ is rule-type data inherited from the Dimensional Atlas.

\subsection{Minimal coupling}

Local-phase invariance $\Psi \to e^{iq\alpha(x)/\hbar}\Psi$ forces the connection $A_\mu \to A_\mu - \partial_\mu \alpha$. The covariant derivative
\begin{equation}
D_\mu = \partial_\mu + \frac{iq}{\hbar} A_\mu
\label{eq:Dmu}
\end{equation}
ensures $D_\mu \Psi$ transforms covariantly. The interacting Klein--Gordon equation
\begin{equation}
\left( D_\mu D^\mu + \frac{m^2 c^2}{\hbar^2} \right) \Psi = 0
\end{equation}
is gauge-covariant. The U(1) gauge structure is therefore \emph{forced} at primitive level by local-phase invariance, not postulated.

\subsection{Conserved current}

Standard manipulation yields the conserved Klein--Gordon current
\begin{equation}
j^\mu = \frac{i\hbar}{2m} \left[ \Psi^* D^\mu \Psi - \Psi (D^\mu \Psi)^* \right],
\end{equation}
satisfying $\partial_\mu j^\mu = 0$ on solutions. The $j^0$ component is sign-indefinite, signalling the well-known negative-probability pathology of Klein--Gordon for matter fields. This pathology is resolved by Stage~R.3's Dirac construction (Section~\ref{sec:r3}).

\subsection{Stage R.1 status}

Stage~R.1 produces:
\begin{itemize}
  \item Lorentz-covariant scalar participation measure (\textsc{forced});
  \item Klein--Gordon equation (\textsc{forced});
  \item Minimal coupling and U(1) gauge structure (\textsc{forced});
  \item Scalar conserved current (\textsc{forced}, sign-indefinite for matter --- superseded by spinor case).
\end{itemize}
Numerical $m$, $q$, $\hbar$, $c$ are inherited.

\section{Arc R Stage R.2 --- Rule-Type Taxonomy and Spin--Statistics}
\label{sec:r2}

Stage~R.2 establishes the structural framework for spin and statistics. We summarise the five-memo sub-program (R.2.1 exchange, R.2.2 Lorentz representations, R.2.3 double-cover topology, R.2.4 Clifford algebra, R.2.5 synthesis) and state the spin--statistics theorem.

\subsection{Exchange dichotomy (R.2.1)}

For two same-type chains $K_A, K_B$, the exchange operation $E_{AB}$ swaps participation labels. Involutive structure ($E_{AB}^2 = \mathrm{id}$) and mutual substitutability of same-type chains restrict the exchange factor $\eta(\tau)$ to
\begin{equation}
\eta(\tau) \in \{+1, -1\}.
\end{equation}
Two rule-type categories emerge: \textbf{Case~P} (bandwidth-sharing-permissive, $\eta = +1$, symmetric joint participation) and \textbf{Case~R} (bandwidth-sharing-restrictive, $\eta = -1$, antisymmetric, vanishing-on-coincidence, Pauli exclusion).

\subsection{Lorentz representations (R.2.2)}

Finite-dimensional representations of the complexified Lorentz algebra $\mathfrak{so}(3,1)_\mathbb{C} \cong \mathfrak{su}(2)_\mathbb{C} \oplus \mathfrak{su}(2)_\mathbb{C}$ are classified by pairs $(j_L, j_R)$ with $j_L, j_R \in \{0, 1/2, 1, 3/2, \dots\}$ and dimension $(2j_L + 1)(2j_R + 1)$. The spin quantum number $s$ runs through $|j_L - j_R| \le s \le j_L + j_R$ via the Pauli--Lubanski Casimir. The exhaustive ladder
\begin{equation}
s \in \{0, 1/2, 1, 3/2, 2, \dots\}
\end{equation}
is rigid in 3+1D by Lie-algebraic theorem.

Integer-spin representations descend to $\mathrm{SO}^+(3,1)$; half-integer representations require the double cover $\mathrm{SL}(2,\mathbb{C})$.

\subsection{Configuration-space topology (R.2.3)}

The fundamental group of the two-identical-chain configuration space $Q_2 = (\mathbb{R}^3 \times \mathbb{R}^3 \setminus \Delta)/S_2$ in 3+1D evaluates to
\begin{equation}
\pi_1(Q_2) = \mathbb{Z}_2.
\label{eq:pi1}
\end{equation}
Centre-of-mass $+$ relative-coordinate decomposition: the relative-coordinate space $(\mathbb{R}^3 \setminus \{0\})/\mathbb{Z}_2 \simeq \mathbb{R}P^2$, with $\pi_1(\mathbb{R}P^2) = \mathbb{Z}_2$.

\textbf{Consequences:}
\begin{enumerate}
  \item Anyonic exchange phases are \emph{forbidden} in 3+1D --- only $\eta \in \{\pm 1\}$ admissible.
  \item The exchange-path generator equals the $2\pi$-rotation generator under the $\mathrm{SO}(3)$-equivariant embedding of the relative coordinate.
\end{enumerate}

\subsection{Clifford algebra uniqueness (R.2.4)}

Half-integer representations require an algebraic frame in which $0$ and $2\pi$ rotations are distinguishable. Consider the structural problem: find the unique real finite-dimensional associative algebra $\mathcal{A}$ generated by elements $e_\mu$ ($\mu = 0,1,2,3$) satisfying:
\begin{itemize}
  \item[(P1)] Lorentz tangent-space compatibility: $\{e_\mu\}$ transforms as a four-vector;
  \item[(P2)] Metric compatibility: symmetric pairings reproduce $\eta_{\mu\nu}$;
  \item[(P3)] Half-integer representation realisation;
  \item[(P4)] Minimal dimension.
\end{itemize}
Pure commutation ($e_\mu e_\nu = e_\nu e_\mu$) fails (P3); pure anticommutation ($e_\mu e_\nu = -e_\nu e_\mu$ strictly) fails (P2). The unique solution is the \textbf{anticommutator} structure
\begin{equation}
\boxed{\{\gamma^\mu, \gamma^\nu\} = 2 \eta^{\mu\nu} \cdot \mathbb{1},}
\label{eq:Cliff}
\end{equation}
defining the real Clifford algebra $\mathrm{Cl}(3,1) \cong M_4(\mathbb{R})$ of dimension~16. The Lorentz generators on the spinor module are
\begin{equation}
\sigma^{\mu\nu} = \frac{i}{2} [\gamma^\mu, \gamma^\nu],
\end{equation}
generating the $\mathrm{SL}(2,\mathbb{C})$ action. Direct computation yields $D(R(2\pi)) = -\mathbb{1}$ on the spinor module --- the rotational double-cover sign emerges automatically from the half-angle in $\exp(-i\theta\sigma^{12}/2)$.

\subsection{Spin--statistics synthesis (R.2.5)}

\begin{theorem}[Spin--Statistics, R1]
For any ED rule-type in 3+1D with Lorentz-covariant internal index structure, the exchange sign $\eta$ and the spin quantum number $s$ are related by
\begin{equation}
\boxed{\eta = (-1)^{2s},}
\label{eq:spinstat}
\end{equation}
with $\eta$ \textsc{forced} to lie in $\{\pm 1\}$ by $\pi_1(Q_2) = \mathbb{Z}_2$ and $s$ \textsc{forced} to lie in $\{0, 1/2, 1, 3/2, 2, \dots\}$ by Lorentz representation theory. Integer-spin rule-types are Case~P (bosonic); half-integer-spin rule-types are Case~R (fermionic).
\end{theorem}

\begin{proof}[Sketch]
The chain runs through six steps:
\begin{enumerate}
  \item Exchange dichotomy $\eta \in \{+1, -1\}$ from rule-type structure (R.2.1).
  \item Topological upgrade $\pi_1(Q_2) = \mathbb{Z}_2$ from spatial dimension~3 + identical-chain exchange (R.2.3).
  \item Geometric theorem: exchange-path generator $=$ $2\pi$-rotation generator (R.2.3).
  \item Lorentz-representation ladder: half-integer requires $\mathrm{SL}(2,\mathbb{C})$ (R.2.2).
  \item Algebraic frame: $\mathrm{Cl}(3,1)$ with $D(R(2\pi)) = -\mathbb{1}$ on spinor module (R.2.4).
  \item Bilinear-pairing closure: Primitive~10 individuation supplies the bilinear $\bar\Psi \Gamma \Psi$ coupling between chain indices, fixing the rule-type-frame coupling (R.2.5).
\end{enumerate}
The minimal-bilinear closure step uses Primitive~10's two-chain individuation pairing made dynamical by Primitive~11. No CANDIDATE assumptions remain; the theorem is \textsc{forced} unconditionally at primitive level.
\end{proof}

\textbf{Dependencies on primitives.} 02, 03, 06, 07, 10, 11, 13.

\subsection{Cl(3,1) uniqueness as a separate theorem}

We isolate the algebraic-frame content of Section~4.4 as a theorem:

\begin{theorem}[Cl(3,1) Uniqueness, R2]
The unique real finite-dimensional associative algebra compatible with (P1) Lorentz tangent-space compatibility, (P2) metric compatibility, (P3) half-integer representation realisation, and (P4) minimal dimension is the real Clifford algebra $\mathrm{Cl}(3,1)$ with $\{\gamma^\mu, \gamma^\nu\} = 2\eta^{\mu\nu} \cdot \mathbb{1}$. Pure commutation fails (P3); pure anticommutation (strict) fails (P2). The anticommutator is structurally mandatory.
\end{theorem}

\begin{proof}[Sketch]
Suppose $e_\mu e_\nu = e_\nu e_\mu$ for all $\mu, \nu$. Then (P2) reduces to $e_\mu e_\nu = (\eta_{\mu\nu}/2)\mathbb{1}$, forcing $e_\mu = 0$ for $\mu \ne \nu$ and rigid scalars on the diagonal. The resulting algebra does not support a half-integer $\mathrm{SL}(2,\mathbb{C})$ representation; (P3) fails.

Suppose $e_\mu e_\nu = -e_\nu e_\mu$ strictly. Then $e_\mu^2 = 0$ for all $\mu$, contradicting (P2) on the diagonal.

The remaining possibility is $e_\mu e_\nu + e_\nu e_\mu = 2\eta_{\mu\nu}\mathbb{1}$, the Clifford-algebra defining relation, which generates $\mathrm{Cl}(3,1)$ uniquely up to isomorphism, with dimension $2^4 = 16$ and basis $\{\mathbb{1}, \gamma^\mu, \gamma^{\mu\nu}, \gamma^{\mu\nu\rho}, \gamma^5\}$.
\end{proof}

\textbf{Dependencies on primitives.} Primitive~06 (Lorentz covariance + metric) plus the half-integer-representation requirement from Theorem~R1.

\section{Arc R Stage R.3 --- Dirac Emergence}
\label{sec:r3}

\subsection{Square-root factorisation}

The Stage~R.1 Klein--Gordon equation applied component-wise to a spinor field $\Psi_\alpha$ does not use the $\gamma^\mu$ generators non-trivially. To use the algebraic frame, a first-order Lorentz-covariant equation built from $\gamma^\mu$, $\partial_\mu$, and a mass scale is required.

Compute:
\begin{align}
(i\gamma^\mu \partial_\mu)(i\gamma^\nu \partial_\nu)
&= -\gamma^\mu \gamma^\nu \partial_\mu \partial_\nu \nonumber \\
&= -\tfrac{1}{2}\{\gamma^\mu, \gamma^\nu\} \partial_\mu \partial_\nu \nonumber \\
&= -\eta^{\mu\nu} \partial_\mu \partial_\nu \cdot \mathbb{1} = -\Box.
\end{align}
Therefore
\begin{equation}
(i\gamma^\mu \partial_\mu - mc/\hbar)(i\gamma^\mu \partial_\mu + mc/\hbar) = -\Box - m^2 c^2/\hbar^2.
\end{equation}
A spinor satisfying the first-order equation
\begin{equation}
\boxed{(i\gamma^\mu \partial_\mu - mc/\hbar)\, \Psi = 0}
\label{eq:Dirac}
\end{equation}
automatically satisfies Klein--Gordon component-wise. This is the \textbf{Dirac equation}.

\subsection{Uniqueness of the first-order form}

Three independent arguments force the first-order form for Case-R rule-types:
\begin{enumerate}
  \item \textbf{Use of the Cl(3,1) structure.} Theorem~R2 establishes $\gamma^\mu$ as primitive-level frame generators. A second-order equation leaves this structure inert. Equation~\eqref{eq:Dirac} is the unique Lorentz-covariant first-order linear operator on $\Psi$ constructible from $\gamma^\mu, \partial_\mu$, and a mass scale.
  \item \textbf{Single time-derivative / positivity.} Klein--Gordon as a second-order spinor equation gives an indefinite inner product. The first-order Dirac form yields a positive-definite current (Section~\ref{sec:current}).
  \item \textbf{Half-integer representation content.} For the $(1/2, 0) \oplus (0, 1/2)$ Dirac module, the minimal Lorentz-covariant first-order equation is unique up to equivalence: it is the Dirac equation.
\end{enumerate}

\begin{theorem}[Dirac Emergence, R3]
Given Stage~R.1 (Klein--Gordon Casimir) and Stage~R.2 (Cl(3,1) frame, half-integer representation content), the unique first-order Lorentz-covariant linear dynamical equation for Case-R rule-type participation measures is
\begin{equation}
(i\gamma^\mu \partial_\mu - mc/\hbar)\, \Psi = 0.
\end{equation}
No additional postulates are required. The equation is structurally \textsc{forced}.
\end{theorem}

\begin{proof}[Sketch]
Lorentz-covariance plus first-order linear structure built from $\gamma^\mu, \partial_\mu, m$ admits exactly the family $a \cdot i\gamma^\mu \partial_\mu \Psi + b \cdot (mc/\hbar)\Psi = 0$ with $a, b$ scalars. Squaring must reproduce Klein--Gordon component-wise (mass-shell consistency); this forces $a = 1, b = -1$ up to overall normalisation. The Cl(3,1) anticommutator $\{\gamma^\mu, \gamma^\nu\} = 2\eta^{\mu\nu}$ ensures that squaring eliminates the cross term $\gamma^\mu \gamma^\nu \partial_\mu \partial_\nu$ via the symmetry of $\partial_\mu \partial_\nu$.
\end{proof}

\textbf{Dependencies on primitives.} All Arc~R stages: R.1 Casimir + R.2 Cl(3,1) frame + R.2 spin--statistics for the Case-R identification.

\subsection{Conserved current}
\label{sec:current}

Multiplying the Dirac equation by $\bar\Psi = \Psi^\dagger \gamma^0$ and combining with the adjoint equation, the mass terms cancel and we obtain $\partial_\mu (\bar\Psi \gamma^\mu \Psi) = 0$, so
\begin{equation}
j^\mu \equiv \bar\Psi \gamma^\mu \Psi
\end{equation}
is conserved. The time component
\begin{equation}
j^0 = \bar\Psi \gamma^0 \Psi = \Psi^\dagger \gamma^0 \gamma^0 \Psi = \Psi^\dagger \Psi = \sum_\alpha |\Psi_\alpha|^2 \ge 0
\end{equation}
is \textbf{positive-definite}, in contrast to the indefinite Klein--Gordon current. The Dirac construction therefore resolves the negative-probability pathology of Klein--Gordon for spin-$1/2$ matter fields. Positivity is structurally tied to the first-order form via the spinor-module structure.

\subsection{Minimal coupling and gauge invariance}

The Stage~R.1 minimal-coupling procedure applies unchanged on the spinor module:
\begin{equation}
(i\gamma^\mu D_\mu - mc/\hbar)\Psi = 0, \qquad D_\mu = \partial_\mu + \frac{iq}{\hbar} A_\mu.
\end{equation}
Under local U(1), $\Psi \to e^{iq\alpha/\hbar}\Psi$, $A_\mu \to A_\mu - \partial_\mu \alpha$, and $D_\mu \Psi \to e^{iq\alpha/\hbar} D_\mu \Psi$. The interacting Dirac equation is gauge-covariant.

\subsection{Non-relativistic reduction and $g = 2$}

In the non-relativistic limit, write $\Psi = (\varphi, \chi)^T \cdot e^{-imc^2 t/\hbar}$ and solve the small-component equation algebraically. Substituting back and keeping leading-order terms yields the \textbf{Pauli equation}:
\begin{equation}
i\hbar \partial_t \varphi = \left[ \frac{(\mathbf{p} - q\mathbf{A})^2}{2m} + qA^0 - \frac{q\hbar}{2m} \boldsymbol{\sigma} \cdot \mathbf{B} \right] \varphi.
\end{equation}
The Zeeman coefficient identifies the gyromagnetic ratio
\begin{equation}
\boxed{g_{\mathrm{Dirac}} = 2,}
\end{equation}
emerging without any parameter tuning from the Cl(3,1) frame and the first-order form. Higher-order QED corrections $g - 2 \approx \alpha/\pi + \dots$ are loop-level content (Arc~Q~\cite{arcq} handles these); the \emph{tree-level} $g = 2$ is a structural prediction of Arc~R.

Dropping spin-dependent terms (projecting onto a spinless participation) yields the Schr\"odinger equation, recovering the Phase-1 non-relativistic limit and closing a consistency loop across the entire QM-emergence program.

\section{Discussion}
\label{sec:discussion}

\subsection{What Arc R forces vs what it leaves open}

\textbf{Arc~R \textsc{forces}, at primitive level:}
\begin{itemize}
  \item Lorentz-covariant scalar participation measure;
  \item Klein--Gordon equation for spin-0 rule-types;
  \item Minimal coupling and U(1) gauge structure from local-phase invariance;
  \item Spin--statistics theorem $\eta = (-1)^{2s}$ (Theorem~R1);
  \item $(j_L, j_R)$ representation classification with exhaustive spin ladder;
  \item Anyon prohibition in 3+1D (from $\pi_1(Q_2) = \mathbb{Z}_2$);
  \item Cl(3,1) frame uniqueness with $\{\gamma^\mu, \gamma^\nu\} = 2\eta^{\mu\nu}$ (Theorem~R2);
  \item $D(R(2\pi)) = -\mathbb{1}$ on the spinor module;
  \item Dirac equation as unique first-order Lorentz-covariant equation on the spinor module (Theorem~R3);
  \item Positive-definite conserved current $j^\mu = \bar\Psi \gamma^\mu \Psi$;
  \item Tree-level Dirac $g = 2$ from non-relativistic reduction;
  \item Phase-1 consistency: Schr\"odinger emerges as the spinless non-relativistic limit.
\end{itemize}

\textbf{Arc~R \textsc{inherits} (does not predict):}
\begin{itemize}
  \item Numerical $\hbar$, $c$;
  \item Per-rule-type mass $m$, charge $q$;
  \item Specific spin assignment per rule-type;
  \item Choice between Dirac, Weyl, and Majorana reality structure for any specific Case-R rule-type;
  \item The number of rule-types occupying any particular structural slot.
\end{itemize}
The form/value separation is preserved throughout. ED's primitive structure determines the \emph{form} of relativistic quantum equations and the \emph{classification} of admissible spin/statistics structures; numerical content is inherited via the Dimensional Atlas.

\subsection{Relationship to Arc M}

Arc~M~\cite{arcm} addresses the chain-mass problem: whether ED primitives constrain the numerical mass values of rule-types. Building on Arc~R's spin--statistics framework and Cl(3,1) frame, Arc~M constructs a candidate bandwidth-signature
\begin{equation}
\sigma_\tau = \hbar \cdot \sqrt{\sum_X w_\tau^X \cdot \langle (\partial_\mu \ln b_\tau^X)(\partial^\mu \ln b_\tau^X) \rangle_\tau}
\end{equation}
as a Lorentz-scalar functional of rule-type bandwidth content. Arc~M closes ``H1-dominant'': $\sigma_\tau$'s Lorentz-scalar form is \textsc{forced}, but mass values, ratios, and hierarchies are \textsc{inherited}. ED fixes the \emph{structure} of mass but not its \emph{values}. Arc~R supplies the foundational framework for Arc~M's analysis.

\subsection{Relationship to Arc Q}

Arc~Q~\cite{arcq} extends ED to the QFT regime: gauge-group structure, interaction vertices, Higgs-like SSB, radiative corrections, generations, second quantisation, and vacuum structure. Three Arc~Q theorems extend the structural framework:
\begin{itemize}
  \item \textbf{GRH unconditional \textsc{forced}} --- at least one massless Case-P $(1/2, 1/2)$ gauge rule-type exists structurally.
  \item \textbf{Canonical (anti-)commutation \textsc{forced}} --- second-quantisation derived directly from Arc~R's spin--statistics + Primitive~10 individuation.
  \item \textbf{UV-FIN \textsc{forced} at primitive level} --- Primitive~01 event-discreteness + Primitive~13 finite intervals + Primitive~04 bounded bandwidth jointly guarantee primitive-level UV finiteness; renormalisation \textsc{refuted} as fundamental.
\end{itemize}
Arc~Q's results presuppose Arc~R's framework: the Cl(3,1) frame supplies the vertex catalogue's Fierz basis; the spin--statistics theorem forces the (anti-)commutation structure; the Dirac equation supplies the matter-sector dynamics for QFT loop calculations.

\subsection{How Arc R sets the stage for the QFT extension}

Arc~R closes the relativistic single-chain quantum sector and supplies the structural prerequisites for Arc~Q:
\begin{itemize}
  \item The Cl(3,1) frame from Theorem~R2 is the algebraic backbone of the Arc~Q vertex catalogue.
  \item Theorem~R1's spin--statistics directly forces Arc~Q's canonical (anti-)commutation relations.
  \item The Dirac equation of Theorem~R3 supplies the matter-sector dynamics for Arc~Q's interaction-vertex analysis.
  \item The minimal-coupling derivation extends to non-Abelian gauge groups without additional structural input.
\end{itemize}
Together, Arc~R + Arc~Q yields a structurally complete, UV-finite quantum field framework whose form is forced and whose numerical content is inherited.

\section{Conclusion}

Arc~R derives relativistic single-chain quantum mechanics as a forced structural consequence of the ED primitive stack. Three structural theorems carry the load: spin--statistics (Theorem~R1), Cl(3,1) uniqueness (Theorem~R2), and Dirac emergence (Theorem~R3). Together they reduce the conventional independent postulates of relativistic quantum theory --- the form of the wave equation, the spin--statistics correspondence, the spinor frame, the gauge-covariant minimal coupling, and the conserved current --- to consequences of a small primitive stack on a 3+1-dimensional event manifold.

The methodological framing --- \emph{form \textsc{forced}, value \textsc{inherited}} --- distinguishes ED's contribution from numerical-prediction frameworks. ED does not predict the numerical value of $\hbar$, the electron mass, or the fine-structure constant. It predicts the \emph{structural form} of the equations these constants appear in: the Klein--Gordon and Dirac equations, the gauge-covariant derivative, the conserved current, the spin ladder, the exchange-statistics dichotomy, and the Cl(3,1) algebraic frame.

Arc~R's tree-level Dirac $g = 2$ emerges from the Cl(3,1) frame and the first-order form without parameter tuning. Its positive-definite conserved current resolves the Klein--Gordon negative-probability pathology for spin-$1/2$ fields. Its anyon prohibition follows from $\pi_1(Q_2) = \mathbb{Z}_2$ --- a topological theorem about 3+1-dimensional configuration space.

These structural results ground the further extensions of Arc~M (chain-mass) and Arc~Q (QFT). Together, Phase-2 of the ED program (Arc~R + Arc~M + Arc~Q) yields a structurally complete, UV-finite quantum field framework. The clean separation between structural prediction and empirical inheritance is preserved at every layer.

Phase-3 directions --- ED to general-relativistic coupling, cosmological structure via the UV-FIN reframing of cosmological-$\Lambda$, possible high-energy empirical signatures near primitive event-discreteness scales --- are unblocked by the Phase-2 closure of which Arc~R is the relativistic foundation.

\begin{thebibliography}{99}

\bibitem{phase1}
A.~Proxmire and Copilot, \emph{Quantum Mechanics as a Structural Consequence of Event-Density Primitives} (Phase-1 closure paper), 2026.

\bibitem{arcm}
A.~Proxmire, \emph{Chain-Mass Structure in Event Density} (Arc~M synthesis), \texttt{chain\_mass\_synthesis.md}, 2026.

\bibitem{arcq}
A.~Proxmire, \emph{QFT Extension of Event Density: Arc~Q Synthesis}, \texttt{arc\_q\_synthesis.md}, 2026.

\bibitem{primitives}
A.~Proxmire, \emph{Event-Density Primitive Stack} (Primitives 01--13), 2026.

\bibitem{phase2syn}
A.~Proxmire, \emph{Phase-2 Global Synthesis}, \texttt{phase2\_synthesis.md}, 2026.

\bibitem{streater}
R.~Streater and A.~Wightman, \emph{PCT, Spin and Statistics, and All That}. Princeton University Press, 1964.

\bibitem{schwinger}
J.~Schwinger, ``On Quantum-Electrodynamics and the Magnetic Moment of the Electron,'' \emph{Phys.\ Rev.} \textbf{73}, 416 (1948).

\bibitem{dirac}
P.~A.~M.~Dirac, ``The Quantum Theory of the Electron,'' \emph{Proc.\ R.\ Soc.\ Lond.\ A} \textbf{117}, 610 (1928).

\bibitem{weyl}
H.~Weyl, \emph{The Theory of Groups and Quantum Mechanics}. Methuen, 1931.

\bibitem{ryder}
L.~H.~Ryder, \emph{Quantum Field Theory}, 2nd ed. Cambridge University Press, 1996.

\end{thebibliography}

\vspace{1em}
\noindent\rule{\linewidth}{0.4pt}
\vspace{0.2em}

\noindent\emph{Manuscript closure: 2026-04-24. Companion documents: Arc~M synthesis~\cite{arcm}, Arc~Q synthesis~\cite{arcq}, Phase-2 synthesis~\cite{phase2syn}. All Arc~R memos available at} \texttt{quantum/foundations/} \emph{in the Event Density repository.}

\end{document}
